\chapter{Гейт условных ожиданий произведения}

Для \(X=\mathbb{RP}^3\), \(C=S^1\) нормированные состояния
\begin{equation}
 \tau_X(f)=\pi^{-2}\int_X f,
 \qquad
 \tau_C(g)=(2\pi)^{-1}\int_C g
\end{equation}
определяют условные ожидания \(E_C=\operatorname{id}\otimes\tau_C\) и
\(E_X=\tau_X\otimes\operatorname{id}\). Для поднятых факторных мод
\begin{equation}
 \frac1{2\pi}\int_K|f(x)|^2=\int_X|f|^2,
 \qquad
 \frac1{\pi^2}\int_K|g(s)|^2=\int_C|g|^2.
\end{equation}
Следовательно, в модуле \(L^2(X)\oplus L^2(C)\)
\begin{equation}
 \|(1_X,1_C)\|^2=\pi^2+2\pi
\end{equation}
без свободных факторных весов.

\begin{proposition}[Conditional-measure pass]
Нормированное интегрирование по spectator-фактору канонически выводит
собственную меру каждого фактора.
\end{proposition}

Однако одна алгебра \(C^\infty(K)\) имеет только одну глобальную
постоянную моду. Если \(F=1_X+1_C=2\) читать одним полем, двойное условное
считывание даёт \(4(\pi^2+2\pi)\). Центрированное разложение устраняет
двойной счёт, но оставляет только один mean-сектор.

\begin{lemma}[Multiplicity obstruction]
Условные ожидания выводят меры уже заданных модулей, но не создают
удвоение постоянной моды. Требуется вывести
\(L^2(X)\oplus L^2(C)\) из конечной алгебры, суперсвязности или
спектральной тройки.
\end{lemma}

На объявленном градуированном модуле
\begin{equation}
 L^2(X)\oplus L^2(C)\oplus\mathbb C^{23}
 \oplus\Omega^1_{\mathrm{int}}(C)
\end{equation}
единая прямая сумма trace--Hodge-форм воспроизводит обе целевые нормы.
Но объявление модуля вручную переносит скрытый выбор из таблицы весов в
таблицу field content.

Машинный аудит сохранён в
\texttt{s2t\_v3\_conditional\_expectation\_results.json}.